\section{Introduction}
\label{sec:introduction}
% Source: docs/RELATED_WORK.md category A, DECISIONS.md s23, s26.1, s26.3, s26.5. Budget: 1 page.

A machine-learning intrusion detector needs labeled malicious traffic to train on, and labeled
malicious traffic is expensive: someone has to attack the network, or wait for someone else to,
and then confirm what happened. Honeypots offer an appealing shortcut. A decoy has no production
purpose, so nothing legitimate should ever reach it; whatever does is unsolicited by definition,
and can be auto-labeled malicious with no human in the loop. Feed that traffic into a detector's
training set and retrain periodically, and the detector adapts to new attack traffic without
anyone hand-labeling a single flow. This design is not hypothetical. Recent systems build
exactly this loop: \citet{matcheswala2024creating} pairs an adaptive honeypot with a network
traffic classifier trained on what the honeypot collects, and \citet{lanka2024intelligent}
applies machine-learning analysis to honeypot-collected data for threat detection. Neither
system asks what happens if the traffic reaching the decoy is not what the auto-labelling policy
assumes it is.

The assumption the whole design rests on is that the label is trustworthy \emph{because} the
source is a decoy. That assumption does not survive contact with an adversary who knows the
policy. The labelling rule — "everything the honeypot sees is malicious" — is public by
construction; there is no way to run the defense without publishing it, since the honeypot has
to be reachable to do its job. An attacker who knows the rule does not need to compromise
anything to write into the defender's training set: sending traffic to the decoy \emph{is}
writing to the training set, with a label the defender chooses to trust unconditionally. The
honeypot is not a passive sensor in this design; it is a write channel the defender opens on
purpose, and the poison budget available through it is not a fraction an attacker has to fight
for — it is the entire honeypot-derived share of the training set, by construction
(Section~\ref{sec:threat-model}).

This paper characterises that channel on a modern, flow-based network intrusion detection loop:
seed data plus honeypot traffic, retrained periodically, scored against held-out production
traffic. We confirm first that the loop does what it claims when the honeypot sees genuine
attacks — the honest case, S0, has to work before an attack on it means anything
(Section~\ref{sec:attack}). We then attack it with A1, benign mimicry: point the same traffic
generator that produces legitimate benign flows at the honeypot instead of production, and let
the auto-labeller do the rest. The result is one mechanism — a label that contradicts the
neighbourhood it sits in — that surfaces as false positives under a fixed detection threshold or
as lost detection under a recalibrated one, not two independent effects; recalibrating relocates
where the damage shows up, it does not remove it. On CICIDS2017 the attack requires
near-duplicate mimicry fidelity and needs no model access, no knowledge of the poison ratio, and
no cost to the attacker beyond producing traffic that already looks like traffic the network
generates for free. We hold every quantitative claim in this paper to a pre-registered,
multiple-comparison-corrected significance bar, and where a pattern in the data does not clear
that bar we report it as unconfirmed rather than round it up to a finding.

We then evaluate defenses. A generic label-consistency check — does this row's stamped label
agree with a model fit on the loop's own data — fails structurally against this attack, because
the poison is not inconsistent with anything the loop itself believes; it agrees with the
defender's own labelling policy, which is exactly what makes it hard to catch from inside the
loop. Cost-of-influence weighting, which we set out to propose as this paper's defense, does not
survive its own pre-registered evaluation: it fails a sensitivity criterion fixed in advance and
does not beat a matched no-skill baseline. What does work, and costs almost nothing to deploy,
is capping the honeypot's effective share of the training set — a one-line change that needs no
label information, no trusted reference data beyond the cap value itself, and no estimate of the
poison ratio, and that no alternative we test dominates. Along the way, measuring the attack on
CICIDS2017 exposes a property of the dataset itself: its benign traffic is narrow and repetitive
enough that an attacker gets most of the way to the damaging fidelity level for free, simply by
replaying the pool's own statistics with no perturbation at all — a fact about this benchmark's
degeneracy, not a general claim about production network traffic.

\subsection*{Contributions}
\begin{enumerate}
  \item A threat model for auto-labelled honeypot data as an attacker-controlled write channel
    into IDS training, characterised — mechanism, fidelity requirement, matched null controls —
    on a modern honeypot-fed retraining loop: the design that older signature-generator attacks
    targeted in a different form \cite{newsome2006paragraph,perdisci2006misleading}, and that
    current honeypot-plus-ML systems reintroduce without evaluating the threat
    (Section~\ref{sec:threat-model}).
  \item The attack characterised as one mechanism, surfaced as either false positives or lost
    detection depending on the threshold policy in force, not two independent channels; damage
    reported as a function of measured mimicry fidelity (twin fraction, not a jitter knob or a
    median distance too coarse to resolve the effect), held to a pre-registered,
    multiple-comparison-corrected significance bar throughout, including a dedicated,
    higher-powered follow-up that confirms one model's detection-loss damage extends well past
    copy-level fidelity, where a larger but lower-powered test could not resolve it
    (Section~\ref{sec:attack}).
  \item The policy-consistency blind spot: why a defense that checks a label's consistency
    against the loop's own data cannot catch poison that agrees with the loop's own labelling
    policy, and why a defense anchored to trusted data outside the loop can
    (Section~\ref{sec:defenses}).
  \item A defense comparison against a proper no-skill baseline: cost-of-influence weighting, our
    original proposal, is a measured negative result that does not beat it; a one-line cap on the
    honeypot's training-set share is undominated by every alternative tested and costs almost
    nothing to run (Section~\ref{sec:defenses}).
  \item A measurement of CICIDS2017's near-degeneracy under standard flow features, including
    that the attack's damage on this dataset traces to near-duplicate, not exact-duplicate,
    benign traffic, and what that implies for which claims this benchmark can and cannot support
    (Section~\ref{sec:degeneracy}).
\end{enumerate}

\begin{figure*}[t]
  \centering
  % F1: loop/threat diagram, drawn in TikZ (vector, no external tooling), not generated from
  % data. Colors match experiments/paper_figures.py's palette (loopblue = BLUE #2A78D6,
  % loopred = RED #E34948, loopmut = MUT #8A8A86, defined in main.tex) so F1 reads as the same
  % document as F2-F6. Source: CLAUDE.md architecture diagram + Section~\ref{sec:threat-model}.
  % Tried single-column (3.33in) first; two labeled inputs, a 4-stage pipeline, and two separate
  % evaluation callouts do not stay legible at that width, so this is figure* (full width).
  \begin{tikzpicture}[
      font=\sffamily\footnotesize,
      node distance=6mm and 7mm,
      box/.style={draw=loopblue, fill=loopbluebg, rounded corners=2pt, minimum height=8mm,
                  text width=20mm, align=center, thick, inner sep=2pt},
      hpbox/.style={box, draw=loopred},
      evalbox/.style={draw=loopmut, fill=white, dashed, rounded corners=2pt, minimum height=8mm,
                       text width=20mm, align=center, inner sep=2pt},
      arr/.style={-{Stealth[length=2mm]}, thick, draw=black!70},
    ]
    % --- inputs ---
    \node[box] (seed) {\texttt{seed\_train}\\[1pt]\lbl{trusted, defender-controlled}};
    \node[hpbox, below=of seed] (pool) {\texttt{honeypot\_pool}\\[1pt]\lbl{attacker sends traffic here}};

    % --- attacker-controlled pipeline ---
    \node[hpbox, right=of pool] (label) {auto-label\\[1pt]``malicious''};
    \node[box, right=of label] (defense) {defense hook\\[1pt]\lbl{optional, e.g. ShareCap}};

    % --- shared pipeline ---
    \node[box, right=16mm of $(label)!0.5!(defense)$, yshift=9mm] (accum) {accumulate};
    \node[box, right=of accum] (retrain) {retrain\\[1pt]\lbl{cold-start each round}};
    \node[box, right=of retrain] (detector) {detector};
    \node[box, right=of detector, text width=24mm] (score) {scores production traffic};

    % --- evaluation, off to the side, never touched by the loop ---
    \node[evalbox, below=10mm of detector] (evalset) {\texttt{trusted\_eval}\\[1pt]\lbl{held out, never written by the loop}};
    \node[evalbox, right=of evalset, text width=24mm] (thresh) {fixed vs.\ recalibrated threshold};

    % --- arrows ---
    \draw[arr] (seed.east) -| ($(seed.east)!0.5!(accum.west)$) |- (accum.west);
    \draw[arr] (pool) -- (label);
    \draw[arr] (label) -- (defense);
    \draw[arr] (defense.east) -| ($(defense.east)!0.5!(accum.west)$) |- (accum.west);
    \draw[arr] (accum) -- (retrain);
    \draw[arr] (retrain) -- (detector);
    \draw[arr] (detector) -- (score);
    \draw[arr] (evalset) -- (thresh);
    \draw[arr, loopmut] (evalset.north) -- (detector.south);
    \draw[arr, loopred, dashed] ($(pool.west)+(-9mm,3mm)$) node[font=\sffamily\scriptsize, loopred, anchor=east]{attacker} -- (pool.west);

    % --- attacker-controlled segment, marked ---
    \begin{pgfonlayer}{background}
      \node[draw=loopred, dashed, thick, rounded corners=4pt, fit=(pool)(label)(defense),
            inner sep=3mm, label={[loopred, font=\sffamily\scriptsize]above:attacker-controlled segment}] {};
    \end{pgfonlayer}
  \end{tikzpicture}
  \caption{The retraining loop. Everything the honeypot logs is auto-labeled malicious by a
  public policy and feeds the same \texttt{accumulate} step as trusted seed data (dashed red
  box: the attacker-controlled segment — the only traffic the attacker chooses is what the
  decoy sees, and that traffic becomes training data unconditionally). \texttt{trusted\_eval} is
  held out and never written by the loop; every result in this paper is scored against it under
  both a fixed and a recalibrated detection threshold.}
  \label{fig:loop-diagram}
\end{figure*}
